\section{Progression-free polynomial sets over finite local rings}
\label{sec:finite-rings}

Asada, Fourakis, Manski, McNew, Miller and Moreland asked in Question~6.5 of
\emph{Subsets of $\mathbb F_q[x]$ free of 3-term geometric progressions}
what changes when the coefficient field is replaced by another finite ring.
We retain their nonunit-ratio convention. The analogue of their degree
ball is the set of polynomials whose coefficients above degree $n$ vanish.
All ratios in this section are polynomials.
\href{https://www.nathanmcnew.com/fqxramsey.pdf}{Author manuscript, pp.~10--12};
\href{https://doi.org/10.1016/j.ffa.2016.10.002}{published paper}.

The Lean connection is an intermediate lemma. In the source,
\texttt{modEq\_orderOf\_of\_div\_pow\_eq} turns a multiplicative equality
into a restriction on exponents, and \texttt{exists\_three\_term\_free}
constructs sets that avoid the resulting additive configurations. An
earlier ProveIt report applied this idea to polynomial valuations over a
field. Over a ring with nilpotents the scalar ratio and its square
act differently. Tracking the prime valuations of the \emph{coefficients
in a basis of the ring} detects forbidden endpoints. The resulting
argument is elementary and independent of the statements about zeros
of $L$-functions.
\href{https://github.com/lsernau/grh-lean/blob/72fead102f663d5c7a407964a13dc0cf9d977207/grh/src/New/Reduction/Data/Existence.lean#L80-L162}{Lean source, lines 80--162};
\href{https://github.com/VladimirReshetnikov/ProveIt/blob/d31039a48d963692e5a5176f47b6603fc9e9d9cf/docs/incoming/grh-progression-densities-and-cubic-prime-race.zip}{earlier field case}.

Let $k=\mathbb F_q$, let $V$ be a $k$-vector space of dimension $b-1$,
and give $R=k\oplus V$ the commutative multiplication
$(u,v)(u',v')=(uu',uv'+u'v)$. Put $Q=|R|=q^b$.
With a fixed basis of $V$, write $f\in R[x]$ as
$(f_0,f_1,\ldots,f_{b-1})$ with entries in $k[x]$.
The ball $B_n$ consists of all such tuples of degrees at most $n$ and has
$Q^{n+1}$ elements. A set is admissible if it contains no triple
$f,rf,r^2f$ with $r\in R[x]$ a nonunit. We work with subsets of
$R[x]\setminus\{0\}$. Retaining zero would change no density.
Write $\Delta(R)$ for the supremum of the upper ball densities of
admissible sets. The natural ball density uses a limit in place of a
limit superior. Reduction of $r=(a,c)$ modulo $V[x]$ detects units:
$r$ is a unit exactly when $a\in k^\times$, as
$(a,c)^{-1}=(a^{-1},-a^{-2}c)$ in this case.

\medskip\noindent\textbf{Theorem (characteristic two).}
Suppose $q$ is even. The set $S$ of nonzero tuples whose monic componentwise
gcd is squarefree is admissible. For every $n\ge2$ it satisfies
\[
 |S\cap B_n|=Q^{n+1}-qQ^{n-1}+q-1.
\]
Its natural density is $1-q/Q^2$. The following bounds include a stronger
construction for upper density:
\begin{equation}
 1-\frac{1}{Q^2}\ \le\ \Delta(R)
 \ \le\ 1-\frac{Q-1}{Q^3-1}
 =1-\frac1{Q^2+Q+1}.
 \label{eq:local-ring-densities}
\end{equation}
For fixed $q$ this determines the leading constant:
$Q^2(1-\Delta(R))\to1$ as $b\to\infty$. All the even-characteristic
assertions hold more generally for a finite commutative local $k$-algebra
$R=k\oplus V$ with nilpotent maximal ideal $V$, provided every $v\in V$
satisfies $v^2=0$. Distinct elements of $V$ may have nonzero product.

\smallskip\noindent\emph{Proof.}
In characteristic two, $r=a+c$ has $r^2=a^2$, since Frobenius
acts coefficientwise on $c\in V[x]$ and all its coefficients square to zero.
In the more general algebra the unit criterion follows by reduction
modulo the nilpotent ideal $V[x]$ and a finite geometric-series inverse. If $a=0$ the last term
of the progression is zero. Otherwise choose a monic irreducible
$P\mid a$. Every component of the last term is divisible by $P^2$.
Conversely every nonzero tuple with a common square factor $P^2$ is a
multiple of the square of the scalar nonunit $P$. Thus $S$ consists
exactly of the nonzero elements divisible by no square of a nonunit.

For any nonzero tuple the indicator that its gcd is squarefree is
$\sum_{d^2\mid\gcd(f_i)}\mu(d)$ over monic $d\in k[x]$. In $B_n$
there are $Q^{n+1-2j}-1$ nonzero multiples of a fixed $d^2$ of degree
$2j\le n$. Unique factorization gives
$\sum_{d\ \mathrm{monic}}\mu(d)z^{\deg d}=1-qz$: it is the reciprocal
of $\sum_{j\ge0}q^j z^j$. Thus only $j=0,1$ contribute after
summing, yielding the count and the asserted natural density.

We next compute the best upper density under the stronger condition that the set
contain no pair $f,a^2f$ with $a\in k[x]$ nonconstant. In a single
ball $B_n$, $n\ge1$, the two layers of highest degrees $n-1,n$ avoid
all such pairs. They have $Q^{n+1}-Q^{n-1}$ elements. For the reverse
bound, map each nonzero $f\in B_n$ to
$x^{2\lfloor(n-\deg f)/2\rfloor}f$ in those two layers. Two distinct
preimages of one point differ by multiplication by a positive even
power of $x$. The map is therefore injective on any pair-avoiding set.

To attain the limiting proportion in a single infinite set, let
$h(f)$ be the largest monic polynomial whose square divides every
component of $f$. In each prime valuation its exponent is half the common
valuation rounded down. Thus
\[
 m(f)=\deg f-2\deg h(f)
\]
is invariant under multiplication by any nonzero scalar polynomial square.
Choose $N_0=-1$ and $N_j=2^{j+1}$ for $j\ge1$. Assign each $f\ne0$
to the unique $j$ with $N_{j-1}<m(f)\le N_j$, and include it exactly
when $\deg f\in\{N_j-1,N_j\}$. All members of a forbidden pair
would have the same assigned $j$ and degrees differing by at least two,
which is impossible. In the top two layers of $B_{N_j}$ a missing
element must have $m(f)\le N_{j-1}$, hence
$\deg h(f)\ge L_j=\lceil(N_j-1-N_{j-1})/2\rceil$.
For any monic $d$ of degree $s\le N_j/2$ the fraction of $B_{N_j}$
divisible componentwise by $d^2$ is $Q^{-2s}$. A union bound
shows that the missed fraction is at most
$\sum_{s\ge L_j}(q/Q^2)^s$, which tends to zero. This proves the
lower bound in \eqref{eq:local-ring-densities} and gives exactly
$1-Q^{-2}$ as the largest attainable upper density under pair avoidance.
The construction also applies for $b=1$.

The natural-density extremum for this same pair condition is also exact.
Each nonzero $f$ factors uniquely as $u h^2$, with $h=h(f)$ monic
and $u$ squarefree in the componentwise gcd sense. For a pair-avoiding
set $E$, let $I_u$ be the set of monic $h$ with $u h^2\in E$.
If $1\in I_u$, there are no other elements. Otherwise
\[
  \sum_{h\in I_u}Q^{-2\deg h}
  \le \sum_{s\ge1}(q/Q^2)^s\le1,
\]
since $q/Q^2\le1/2$. In either case the mass is at most one.
Multiplying by $Q^{-\deg u}$ and summing over base tuples up to degree $n$
gives
\[
  \sum_{f\in E\cap B_n}Q^{-\deg f}
  \ \le\ \sum_{u\in S\cap B_n}Q^{-\deg u}.
\]
The number of all tuples of exact degree $t$ is $(Q-1)Q^t$.
Whenever a set has natural ball density $\delta$, its exact-degree
proportions tend to $\delta$ as well, by subtracting consecutive ball
counts. Division of this last inequality by $(n+1)(Q-1)$ and
convergence of arithmetic means now bound the natural density of $E$ by
$1-q/Q^2$. The squarefree set $S$ attains it. In particular the two
extrema for pair avoidance differ by exactly $(q-1)/Q^2$.

For the upper bound for triple avoidance take any fixed scalar monic linear $P$. Every
nonzero tuple has an exact common $P$-valuation $v$, the minimum of
the component valuations with $v_P(0)=+\infty$. Partition each chain
$h,Ph,P^2h,\ldots$ with $v(h)=0$ into triples starting at exponents
$0,3,6,\ldots$. Multiplication by $P$ is injective and increases the
maximum component degree by one. Hence the triples are disjoint, and
every triple lying within $B_n$ forces at least one deletion. Their
number is exactly
\[
 T_n=\sum_{0\le 3j\le n-2}
 (Q^{n-1-3j}-Q^{n-2-3j}).
\]
Indeed their first terms have degree at most $n-2$ and valuation $3j$.
Dividing by $Q^{n+1}$ and summing the limiting geometric series gives
$(1-Q^{-1})Q^{-2}/(1-Q^{-3})$. This argument works in every
characteristic, including $b=1$. Equation~\eqref{eq:local-ring-densities}
now squeezes $Q^2(1-\Delta(R))$ between $Q^2/(Q^2+Q+1)$ and $1$.\hfill$\square$

\medskip\noindent\textbf{Theorem (a strict gap for triple avoidance).}
Let $D(R)$ be the supremum of natural densities of admissible sets
whose ball densities converge. In the characteristic-two family,
\[
 1-\frac q{Q^2}\le D(R)
 \le 1-\frac q{Q^2}(1-Q^{-1})^q.
\]
Thus as $b\to\infty$ with $q$ fixed, $Q^2(1-D(R))\to q$.
For every $b\ge2$ in this family, $D(R)<\Delta(R)$.

\smallskip\noindent\emph{Proof.}
The lower bound comes from the squarefree construction. To bound natural
density from above, let $C$ consist of nonzero tuples whose gcd has no
linear factor. By reduction modulo the product of the $q$ monic linears,
the natural density of $C$ is $(1-Q^{-1})^q$. For each $g\in C$ form
the $q$ triples $g,Pg,P^2g$, one for each monic linear $P$. Two different
centers give disjoint collections: the common linear valuations of a
member determine the multiplying polynomial $1$, $P$ or $P^2$ uniquely.
Within one collection only the center $g$ is shared. Give a tuple $f$
of degree $d$ the weight $Q^{-d}$. Omitting $g$ costs $Q^{-\deg g}$.
If $g$ is present, at least one of $Pg,P^2g$ is missing for each $P$,
which costs at least $qQ^{-2-\deg g}$. In either case the omitted
weight is at least $qQ^{-2-\deg g}$. Sum over centers in $B_{n-2}$,
so every tuple counted is in $B_n$. Divide by $(n+1)(Q-1)$ and pass to
the limit using the exact-degree proportions of both $C$ and the
omitted set as above. This gives the upper bound. For $b\ge2$ one has
$q(1-Q^{-1})^q>q(1-q/Q)\ge q-1\ge1$, by strict Bernoulli inequality.
Consequently this natural-density upper bound is strictly below the
upper-density construction $1-Q^{-2}$.\hfill$\square$

For $R=\mathbb F_2[\epsilon]/(\epsilon^2)$, the construction has
natural density $7/8$. Every admissible set with a natural density has
density at most $119/128$, while the block construction has upper density
$15/16$. Every admissible subset of $\mathbb F_2[x]$ has
upper density at most $6/7$. More generally any even $q$ and $b\ge2$
give a lower bound at least $1-q^{-3}$, which exceeds the residue-field
upper bound $1-1/(q^2+q+1)$ because $q^3>q^2+q+1$.
The same comparison for an arbitrary $q$ holds in square-zero
extensions of dimension $b\ge3$, by the following calculation.

\medskip\noindent\textbf{Theorem (odd characteristic).}
If $q$ is odd, the nonzero elements divisible by no square of a nonunit
form an admissible set of natural density
\begin{equation}
 \prod_{P\ \mathrm{monic\ irreducible}}
 \bigl(1-z_P^{b+1}+z_P^{b+2}-z_P^{2b+1}\bigr)
 \ \ge\ 1-q^{-b},
 \quad z_P=q^{-\deg P}.
 \label{eq:local-ring-odd}
\end{equation}

\smallskip\noindent\emph{Proof.}
Every forbidden triple ends at a multiple of a nonunit square. Discard
the tuples with $f_0=0$ temporarily, a set of proportion $q^{-(n+1)}$.
We claim the nonzero scalar part $f_0$ of an endpoint forces at least one
prime $P$ to satisfy one of the following two mutually exclusive conditions:
\[
\begin{array}{ll}
 v_P(f_0)=2, &P\mid f_i\quad(1\le i<b),\\
 v_P(f_0)\ge3, &P^2\mid f_i\quad(1\le i<b).
\end{array}
\]
Conversely either condition suffices for square divisibility. To see
necessity write $f=(a,c)^2(h_0,h)=(a^2h_0,a^2h+2ach_0)$
and take $P\mid a$. All nilpotent coordinates are multiples of $P$.
If $v_P(a^2 h_0)\ge3$, either $P^2\mid a$ or $P\mid h_0$,
and all are multiples of $P^2$. To prove sufficiency in the second
case divide by the pure square $P^2$. In the first case set
$h_0=f_0/P^2$. Since both $2$ and $h_0$ are invertible modulo $P$,
choose $c_i$ solving $2c_i h_0\equiv f_i/P\pmod P$. Then
$h_i=(f_i-2Pc_i h_0)/P^2$ is a polynomial for each $i$ and
$f=(P,c)^2(h_0,h)$.

For fixed $P$, uniform residues of independent polynomial coordinates
modulo $P^3$ give total local obstruction probability
$(z_P^2-z_P^3)z_P^{b-1}+z_P^3 z_P^{2(b-1)}$.
The Chinese remainder theorem gives independence for any finite set of
primes once $n$ is sufficiently large. For nonzero $f_0$ an obstruction
with $\deg P=d$ requires $2d\le n$, and has probability at most
$q^{-(b+1)d}$ because it needs $P^2\mid f_0$ and $P\mid f_i$ for all
$i>0$. Thus the omitted large primes have total probability at most
$\sum_{d>D} q^{-bd}$, tending to zero uniformly in $n$.
This gives the product. Since $2b+1\ge b+2$ for $b\ge1$, every
factor is at least $1-z_P^{b+1}$. Their product is $1-q^{-b}$ by
unique factorization.\hfill$\square$

These bounds answer Question~6.5 for an infinite family of finite local
rings. They give both density extrema for endpoint-pair avoidance in characteristic two, and the even
case determines the leading asymptotic of the optimal deficit. The
exact triple-avoidance optimum even for the dual numbers over
$\mathbb F_2$, and comparable results for general finite coefficient
rings remain undetermined here.
